\documentclass[10pt,a4paper]{amsart}
\usepackage{amsmath,amssymb,amsthm,graphicx}
\usepackage{verbatim}
\newenvironment{reference}{\comment}{\endcomment}
%\newenvironment{reference}{}{}
\newenvironment{slogan}{\comment}{\endcomment}
\newenvironment{history}{\comment}{\endcomment}

% For commutative diagrams we use Xy-pic
\usepackage[all]{xy}

% We use 2cell for 2-commutative diagrams.
\xyoption{2cell}
% Theorem environments.
%
\theoremstyle{plain}
\newtheorem{theorem}[subsection]{Theorem}
\newtheorem{proposition}[subsection]{Proposition}
\newtheorem{lemma}[subsection]{Lemma}

\theoremstyle{definition}
\newtheorem{definition}[subsection]{Definition}
\newtheorem{example}[subsection]{Example}
\newtheorem{exercise}[subsection]{Exercise}
\newtheorem{situation}[subsection]{Situation}

\theoremstyle{remark}
\newtheorem{remark}[subsection]{Remark}
\newtheorem{remarks}[subsection]{Remarks}

\numberwithin{equation}{subsection}

% Macros
%
\def\lim{\mathop{\mathrm{lim}}\nolimits}
\def\colim{\mathop{\mathrm{colim}}\nolimits}
\def\Spec{\mathop{\mathrm{Spec}}}
\def\Hom{\mathop{\mathrm{Hom}}\nolimits}
\def\Ext{\mathop{\mathrm{Ext}}\nolimits}
\def\SheafHom{\mathop{\mathcal{H}\!\mathit{om}}\nolimits}
\def\SheafExt{\mathop{\mathcal{E}\!\mathit{xt}}\nolimits}
\def\Sch{\mathit{Sch}}
\def\Mor{\mathop{\mathrm{Mor}}\nolimits}
\def\Ob{\mathop{\mathrm{Ob}}\nolimits}
\def\Sh{\mathop{\mathit{Sh}}\nolimits}
\def\NL{\mathop{N\!L}\nolimits}
\def\CH{\mathop{\mathrm{CH}}\nolimits}
\def\proetale{{pro\text{-}\acute{e}tale}}
\def\etale{{\acute{e}tale}}
\def\QCoh{\mathit{QCoh}}
\def\Ker{\mathop{\mathrm{Ker}}}
\def\Im{\mathop{\mathrm{Im}}}
\def\Coker{\mathop{\mathrm{Coker}}}
\def\Coim{\mathop{\mathrm{Coim}}}

% Boxtimes
%
\DeclareMathSymbol{\boxtimes}{\mathbin}{AMSa}{"02}

%
% Macros for moduli stacks/spaces
%
\def\QCohstack{\mathcal{QC}\!\mathit{oh}}
\def\Cohstack{\mathcal{C}\!\mathit{oh}}
\def\Spacesstack{\mathcal{S}\!\mathit{paces}}
\def\Quotfunctor{\mathrm{Quot}}
\def\Hilbfunctor{\mathrm{Hilb}}
\def\Curvesstack{\mathcal{C}\!\mathit{urves}}
\def\Polarizedstack{\mathcal{P}\!\mathit{olarized}}
\def\Complexesstack{\mathcal{C}\!\mathit{omplexes}}
% \Pic is the operator that assigns to X its picard group, usage \Pic(X)
% \Picardstack_{X/B} denotes the Picard stack of X over B
% \Picardfunctor_{X/B} denotes the Picard functor of X over B
\def\Pic{\mathop{\mathrm{Pic}}\nolimits}
\def\Picardstack{\mathcal{P}\!\mathit{ic}}
\def\Picardfunctor{\mathrm{Pic}}
\def\Deformationcategory{\mathcal{D}\!\mathit{ef}}
\usepackage{cleveref}
\crefname{originallemma}{original lemma}{original lemmas}
\Crefname{originallemma}{Original lemma}{Original lemmas}
\makeatletter
\newlabel{algebra-lemma-total-ring-fractions-no-embedded-points}{{02LX}{source}{}{sourcepointer-1}{}}
\newlabel{algebra-lemma-total-ring-fractions-no-embedded-points@cref}{{[originallemma][0][]02LX}{[1][0][]source}{}{}{}}
\makeatother
\crefname{originalsection}{original section}{original sections}
\Crefname{originalsection}{Original section}{Original sections}
\makeatletter
\newlabel{decent-spaces-section-residue-fields-henselian-local-rings}{{0EMV}{source}{}{sourcepointer-2}{}}
\newlabel{decent-spaces-section-residue-fields-henselian-local-rings@cref}{{[originalsection][0][]0EMV}{[1][0][]source}{}{}{}}
\makeatother
\crefname{originallemma}{original lemma}{original lemmas}
\Crefname{originallemma}{Original lemma}{Original lemmas}
\makeatletter
\newlabel{more-algebra-lemma-dumb-properties-henselization}{{07QM}{source}{}{sourcepointer-3}{}}
\newlabel{more-algebra-lemma-dumb-properties-henselization@cref}{{[originallemma][0][]07QM}{[1][0][]source}{}{}{}}
\makeatother
\crefname{originallemma}{original lemma}{original lemmas}
\Crefname{originallemma}{Original lemma}{Original lemmas}
\makeatletter
\newlabel{more-algebra-lemma-henselization-dimension}{{06LK}{source}{}{sourcepointer-4}{}}
\newlabel{more-algebra-lemma-henselization-dimension@cref}{{[originallemma][0][]06LK}{[1][0][]source}{}{}{}}
\makeatother
\crefname{originallemma}{original lemma}{original lemmas}
\Crefname{originallemma}{Original lemma}{Original lemmas}
\makeatletter
\newlabel{more-algebra-lemma-henselization-noetherian}{{06LJ}{source}{}{sourcepointer-5}{}}
\newlabel{more-algebra-lemma-henselization-noetherian@cref}{{[originallemma][0][]06LJ}{[1][0][]source}{}{}{}}
\makeatother
\crefname{originallemma}{original lemma}{original lemmas}
\Crefname{originallemma}{Original lemma}{Original lemmas}
\makeatletter
\newlabel{more-algebra-lemma-henselization-reduced}{{06DH}{source}{}{sourcepointer-6}{}}
\newlabel{more-algebra-lemma-henselization-reduced@cref}{{[originallemma][0][]06DH}{[1][0][]source}{}{}{}}
\makeatother
\crefname{originallemma}{original lemma}{original lemmas}
\Crefname{originallemma}{Original lemma}{Original lemmas}
\makeatletter
\newlabel{morphisms-lemma-integral-scheme-rational-functions}{{01RV}{source}{}{sourcepointer-7}{}}
\newlabel{morphisms-lemma-integral-scheme-rational-functions@cref}{{[originallemma][0][]01RV}{[1][0][]source}{}{}{}}
\makeatother
\crefname{originalsection}{original section}{original sections}
\Crefname{originalsection}{Original section}{Original sections}
\makeatletter
\newlabel{spaces-morphisms-section-rational-maps}{{0EML}{source}{}{sourcepointer-8}{}}
\newlabel{spaces-morphisms-section-rational-maps@cref}{{[originalsection][0][]0EML}{[1][0][]source}{}{}{}}
\makeatother
\title{The henselian local ring at the generic point of a prime divisor is a reduced one dimensional local ring receiving the function field}
\author{The Stacks Project Authors}
\date{stacks-project:spaces-over-fields.tex:local-snapshot-20261012; GNU FDL 1.2 or later}
\begin{document}
\maketitle
\noindent\textit{Editorial note.} All delivered mathematical text is the author's own, copied from the pinned Stacks Project chapter file spaces-over-fields.tex; no statement or proof step was rewritten, reordered or removed.\par\medskip
\noindent\textit{Editorial note.} The counted proof is the author's complete argument: the reduction to an elementary etale neighbourhood in which the local ring is one dimensional, reduced and Noetherian, the composition of rational maps producing an injective map of function fields with the identification of the total quotient ring, and the passage to the henselization using its properties and faithful flatness.\par\medskip
\noindent\textit{Editorial note.} The label lemma-order-vanishing names the counted statement, and the proof environment immediately following it in the source is delivered verbatim as its proof.\par\medskip
\noindent\textit{Editorial note.} The counted statement is carried by the author's own lemma environment in the chapter source.\par\medskip
\noindent\textit{Editorial note.} Pre-open gate: python operations/agent-authoring/tools/precard\_gate.py --source spaces-over-fields.tex --label lemma-order-vanishing --start 779 --end 835 --local-id 28061 returned RESULT=OK (L28061 free, occupancy FREE, public interval clear). Candidate taken from sources/stacks-spaces-over-fields/candidates-halfA.json.\par\medskip
\noindent\textit{Editorial note.} The proof ends with the sentence 'We omit the verification that the map is independent of the choice of (U, u) ...; a slightly better approach would be to first observe that colim Q(O\_\{U,u\}) = Q(O\_\{X,xi\}\textasciicircum{}h\}).' — the omitted verification is the independence of a map already constructed from the choices made, a formality for which the author also indicates the cleaner route, while the ring theoretic content (reducedness, dimension, Noetherianness and injectivity) is proved in full, so the statement is counted and the sentence is recorded here verbatim.\par\medskip
\noindent\textit{Editorial note.} The statement and proof are copied from the pinned chapter file; the cited inputs are kept as external pointers with their canonical Stacks tags.\par\medskip
\noindent\textit{Allowed standard prerequisites (editorial).}\begin{itemize}
\item Elementary etale neighbourhoods of a point. Decent Spaces, the section on residue fields and henselian local rings, in particular the existence of an elementary etale neighbourhood and the identification of the henselian local ring as the henselization of the local ring of that neighbourhood. Supplies the elementary etale neighbourhood (U, u) of (X, xi) and the identification of the henselian local ring as the henselization of the local ring of U at u.
\item Function fields of integral schemes and total rings of fractions. Morphisms, the lemma describing the function field of an integral scheme, and Algebra, the lemma identifying the total ring of fractions for a reduced ring whose components all pass through a point. Supply the identification of the function field of U with the total quotient ring of the local ring at u under the choice of U.
\item Properties of the henselization of a local ring. More on Algebra, the lemmas that the henselization of a reduced ring is reduced, that it preserves dimension and Noetherianness, and that the local ring maps faithfully flatly into its henselization. Supply the reducedness, dimension and Noetherianness of the henselian local ring and the injectivity of the map on total rings of fractions.
\end{itemize}\medskip
\ifdefined\hypertarget\hypertarget{sourcepointer-1}{}\fi
\noindent\textit{Source reference.} Original lemma 02LX: sources/stacks-chow/tags.txt, source lines 3424--3424. Prerequisite pointer invoked at proof line(s) 812; it names a standard fact of a sibling chapter that the counted argument uses, and every step of the counted proof built on it is retained verbatim in the delivered body.\par\medskip
\ifdefined\hypertarget\hypertarget{sourcepointer-2}{}\fi
\noindent\textit{Source reference.} Original section 0EMV: sources/stacks-chow/tags.txt, source lines 18320--18320. Prerequisite pointer invoked at proof line(s) 794; it names a standard fact of a sibling chapter that the counted argument uses, and every step of the counted proof built on it is retained verbatim in the delivered body.\par\medskip
\ifdefined\hypertarget\hypertarget{sourcepointer-3}{}\fi
\noindent\textit{Source reference.} Original lemma 07QM: sources/stacks-chow/tags.txt, source lines 9837--9837. Prerequisite pointer invoked at proof line(s) 826; it names a standard fact of a sibling chapter that the counted argument uses, and every step of the counted proof built on it is retained verbatim in the delivered body.\par\medskip
\ifdefined\hypertarget\hypertarget{sourcepointer-4}{}\fi
\noindent\textit{Source reference.} Original lemma 06LK: sources/stacks-chow/tags.txt, source lines 8469--8469. Prerequisite pointer invoked at proof line(s) 822; it names a standard fact of a sibling chapter that the counted argument uses, and every step of the counted proof built on it is retained verbatim in the delivered body.\par\medskip
\ifdefined\hypertarget\hypertarget{sourcepointer-5}{}\fi
\noindent\textit{Source reference.} Original lemma 06LJ: sources/stacks-chow/tags.txt, source lines 8468--8468. Prerequisite pointer invoked at proof line(s) 823; it names a standard fact of a sibling chapter that the counted argument uses, and every step of the counted proof built on it is retained verbatim in the delivered body.\par\medskip
\ifdefined\hypertarget\hypertarget{sourcepointer-6}{}\fi
\noindent\textit{Source reference.} Original lemma 06DH: sources/stacks-chow/tags.txt, source lines 8181--8181. Prerequisite pointer invoked at proof line(s) 821; it names a standard fact of a sibling chapter that the counted argument uses, and every step of the counted proof built on it is retained verbatim in the delivered body.\par\medskip
\ifdefined\hypertarget\hypertarget{sourcepointer-7}{}\fi
\noindent\textit{Source reference.} Original lemma 01RV: sources/stacks-chow/tags.txt, source lines 2316--2316. Prerequisite pointer invoked at proof line(s) 810; it names a standard fact of a sibling chapter that the counted argument uses, and every step of the counted proof built on it is retained verbatim in the delivered body.\par\medskip
\ifdefined\hypertarget\hypertarget{sourcepointer-8}{}\fi
\noindent\textit{Source reference.} Original section 0EML: sources/stacks-chow/tags.txt, source lines 18311--18311. Prerequisite pointer invoked at proof line(s) 806; it names a standard fact of a sibling chapter that the counted argument uses, and every step of the counted proof built on it is retained verbatim in the delivered body.\par\medskip
\begin{lemma}
\label{lemma-order-vanishing}
Let $S$ be a scheme. Let $X$ be a locally Noetherian integral algebraic space
over $S$. Let $Z \subset X$ be a prime divisor and let $\xi \in |Z|$ be
the generic point. Then the henselian local ring $\mathcal{O}_{X, \xi}^h$
is a reduced $1$-dimensional Noetherian local ring and
there is a canonical injective map
$$
R(X) \longrightarrow Q(\mathcal{O}_{X, \xi}^h)
$$
from the function field $R(X)$ of $X$ into the total ring of fractions.
\end{lemma}

\begin{proof}
We will use the results of Decent Spaces, Section
\ref{decent-spaces-section-residue-fields-henselian-local-rings}.
Let $(U, u) \to (X, \xi)$ be an elementary \'etale neighbourhood.
Observe that $U$ is locally Noetherian and reduced.
Thus $\mathcal{O}_{U, u}$ is a $1$-dimensional
(by our definition of prime divisors)
reduced Noetherian ring.
After replacing $U$ by an affine open neighbourhood of $u$
we may assume $U$ is Noetherian and affine.
After replacing $U$ by a smaller open, we may assume every irreducible
component of $U$ passes through $u$.
Since $U \to X$ is open and $X$ irreducible, $U \to X$ is dominant.
Hence we obtain a ring map $R(X) \to R(U)$ by composing rational maps, see
Morphisms of Spaces, Section \ref{spaces-morphisms-section-rational-maps}.
Since $R(X)$ is a field, this map is injective.
By our choice of $U$ we see that $R(U)$ is the total quotient
ring $Q(\mathcal{O}_{U, u})$, see
Morphisms, Lemma \ref{morphisms-lemma-integral-scheme-rational-functions}
and
Algebra, Lemma \ref{algebra-lemma-total-ring-fractions-no-embedded-points}.

\medskip\noindent
At this point we have proved all the statements in the lemma with
$\mathcal{O}_{U, u}$ in stead of $\mathcal{O}_{X, \xi}^h$.
However, $\mathcal{O}_{X, \xi}^h$ is the henselization of
$\mathcal{O}_{U, u}$. Thus $\mathcal{O}_{X, \xi}^h$ is a $1$-dimensional
reduced Noetherian ring, see
More on Algebra, Lemmas
\ref{more-algebra-lemma-henselization-reduced},
\ref{more-algebra-lemma-henselization-dimension}, and
\ref{more-algebra-lemma-henselization-noetherian}.
Since $\mathcal{O}_{U, u} \to \mathcal{O}_{X, \xi}^h$ is
faithfully flat by More on Algebra, Lemma
\ref{more-algebra-lemma-dumb-properties-henselization}
it sends nonzerodivisors to nonzerodivisors.
Therefore we obtain a canonical map
$Q(\mathcal{O}_{U, u}) \to Q(\mathcal{O}_{X, \xi}^h)$
and we obtain our map.
We omit the verification that the map is independent of
the choice of $(U, u) \to (X, x)$; a slightly better
approach would be to first observe that
$\colim Q(\mathcal{O}_{U, u}) = Q(\mathcal{O}_{X, \xi}^h)$.
\end{proof}

\end{document}
